\section{Computational Resources and Project Management}
\label{sec:app1_resources_management}

\subsection{Conception of Novel Solutions}
We started from the isotropy literature \citep{ethayarajh-2019-contextual,rudman-etal-2022-isoscore}, which studies English text with older encoders. Two gaps followed: no study covers the scripts and tokenization levels of Indian languages, and no study compares a contrastive embedding model with a generative model of the same family. Early runs showed a sharp change in every metric at the final norm, so we added the pre/post-norm split and the norm breakdown to the method.

\subsection{Computational Environment}
Table~\ref{tab:app1_compute_resources} lists the resources used. 

\begin{table}[t]
\centering
\small
\begin{tabular}{l>{\raggedright\arraybackslash}p{3.6cm}}
\toprule
Resource Component & Specification \\
\midrule
Hardware Accelerator & NVIDIA Tesla T4 (Kaggle T4$\times$2 allocation) \\
GPU Video Memory & 16\,GB per GPU; float32 model weights require approximately 1.2--4.0\,GB \\
System Memory (RAM) & 29\,GB \\
CPU Cores & 4 \\
Persistent Storage & 20\,GB of saved output space in \texttt{/kaggle/working}; approximately 130\,MB of results per model, plus model weights and dataset cache \\
Software Stack & Python; PyTorch, Transformers.  \\
\bottomrule
\end{tabular}
\caption{Computational resources and environment for Approach 1.}
\label{tab:app1_compute_resources}
\end{table}

